\documentclass{article}
\usepackage[utf8]{inputenc}
\usepackage{amsmath}
\usepackage{amssymb}

\title{LaTeX 编辑器示例文档}
\author{Code Reader}
\date{\today}

\begin{document}

\maketitle

\section{简介}

这是一个用于测试 LaTeX 编辑器功能的示例文档。本文档包含了常见的 LaTeX 元素，包括文本格式、数学公式、列表和环境等。

\section{文本格式}

LaTeX 支持多种文本格式：

\begin{itemize}
    \item \textbf{粗体文本} 使用 \texttt{\textbackslash textbf} 命令
    \item \textit{斜体文本} 使用 \texttt{\textbackslash textit} 命令
    \item \underline{下划线文本} 使用 \texttt{\textbackslash underline} 命令
    \item \texttt{等宽字体} 使用 \texttt{\textbackslash texttt} 命令
\end{itemize}

\section{数学公式}

\subsection{行内公式}

行内公式可以使用美元符号包围，例如 $E = mc^2$ 是爱因斯坦著名的质能方程。勾股定理可以表示为 $a^2 + b^2 = c^2$。

\subsection{块级公式}

著名的欧拉公式：
$$e^{i\pi} + 1 = 0$$

二次方程的求根公式：
\begin{equation}
x = \frac{-b \pm \sqrt{b^2 - 4ac}}{2a}
\end{equation}

多行对齐公式：
\begin{align}
(a + b)^2 &= a^2 + 2ab + b^2 \\
(a - b)^2 &= a^2 - 2ab + b^2 \\
(a + b)(a - b) &= a^2 - b^2
\end{align}

\subsection{复杂数学表达式}

矩阵表示：
$$
\mathbf{A} = \begin{pmatrix}
a_{11} & a_{12} & a_{13} \\
a_{21} & a_{22} & a_{23} \\
a_{31} & a_{32} & a_{33}
\end{pmatrix}
$$

积分公式：
$$\int_{0}^{\infty} e^{-x^2} dx = \frac{\sqrt{\pi}}{2}$$

求和公式：
$$\sum_{n=1}^{\infty} \frac{1}{n^2} = \frac{\pi^2}{6}$$

\section{列表}

\subsection{无序列表}

\begin{itemize}
    \item 第一项
    \item 第二项
    \item 第三项
\end{itemize}

\subsection{有序列表}

\begin{enumerate}
    \item 步骤一：准备环境
    \item 步骤二：编写代码
    \item 步骤三：测试验证
    \item 步骤四：部署上线
\end{enumerate}

\section{特殊字符}

LaTeX 中的特殊字符需要转义：

\begin{itemize}
    \item 百分号：\%
    \item 美元符号：\$
    \item 与符号：\&
    \item 井号：\#
    \item 下划线：\_
    \item 花括号：\{ \}
\end{itemize}

\section{希腊字母}

常用希腊字母：$\alpha$, $\beta$, $\gamma$, $\delta$, $\epsilon$, $\theta$, $\lambda$, $\mu$, $\pi$, $\sigma$, $\phi$, $\omega$

大写希腊字母：$\Gamma$, $\Delta$, $\Theta$, $\Lambda$, $\Pi$, $\Sigma$, $\Phi$, $\Omega$

\section{结论}

本文档展示了 LaTeX 编辑器支持的主要功能。通过这些示例，可以验证编辑器的解析和渲染能力是否正常工作。

\end{document}
